% FILE: appendix/OC_1_3_3_VERDICT_INVARIANT_MINIMALITY_FULL_PROOF.tex
% OC Core 1.3.3 scientific closure artifact.

\section{OC 1.3.3 Verdict-Invariant Minimality Route}
\label{app:oc133-minimality}

\paragraph{Status.}
This appendix is not a public theorem proof in Core 1.3.3.  The public claim
ledger demotes \(T133\)-MIN to a witness-replay route.  The route records what
would be needed for a future minimality theorem; it does not promote global or
verdict-invariant tuple minimality as proved.

\paragraph{Route claim.}
For the OC verdict class used in Core 1.3.3, removing any required component
from the tuple is expected to change at least one admissibility, liveness,
collapse, identity, or release-governance verdict.  In this release that
expectation is maintained as a checked route, not as a theorem.

\paragraph{Witness method.}
Each component has a witness pair \((K,K')\) that agrees on all other
components and differs on one verdict only because the removed component is no
longer observable.

\paragraph{Witness replay.}
The witness matrix in `proofs/minimality/WITNESS_PAIRS.json` supplies pairs for
\(\Omega,A,P,J,\Theta,\partial\Omega,C,k\), carrier, residue, and morphism
support.  The executable replay checks that each row changes a declared verdict
when the target component is removed while the other component slots are marked
as shared.  This is route evidence.  It is insufficient by itself for a public
minimality theorem because the full formal languages, forgetful maps, and
verdict function remain future work.

\paragraph{Demotion boundary.}
Any sentence reading this appendix as a full proof of tuple minimality is
superseded by the claim ledger and theorem registry for Core 1.3.3.  The
release keeps the route because it is useful to reviewers, but it blocks public
promotion of the minimality theorem until stronger formal artifacts exist.
